\documentclass[11pt]{article}
\usepackage[margin=1in]{geometry}
\usepackage{amsmath,amssymb,amsthm,mathtools}
\usepackage[T1]{fontenc}
\usepackage{lmodern,microtype}
\usepackage[colorlinks=true,linkcolor=blue,citecolor=blue,urlcolor=blue]{hyperref}
\newtheorem{theorem}{Theorem}
\newtheorem{lemma}[theorem]{Lemma}
\newtheorem{proposition}[theorem]{Proposition}
\newcommand{\norm}[1]{\left\|#1\right\|}
\newcommand{\Tt}{\mathbb T_2}
\title{Failure of local controllability at time $21/10$\newline for a shallow-water tank}
\author{Kaiwen Jin\\\small\texttt{yui.ui.siki@gmail.com}}
\date{October 12, 2026}
\begin{document}
\maketitle
\begin{abstract}
We consider a one-dimensional tank governed by the normalized Saint--Venant
equations, with the tank acceleration as its only control. We show that local
exact controllability around the flat equilibrium, with uniformly small
classical trajectories in $C^1$, fails at time $T=21/10$. The terminal states
form an explicit family tending to equilibrium, and the tank starts and ends
with zero velocity and displacement. Thus the sufficiency of every time
$T>2$ is false. The proof computes the full quadratic endpoint in a spatial
primitive, separates it uniformly from a prescribed terminal profile, and
controls the nonlinear error in the same supremum norm. The least time for
local controllability is not determined here.
\end{abstract}

\section{The problem and the counterexample}
The motion of a shallow inviscid fluid in a horizontally accelerated tank is
described by the Saint--Venant equations. After normalizing the tank length,
gravity and equilibrium depth to one, the system is
\begin{equation}\label{eq:tank}
 H_t+(Hv)_x=0,\qquad
 v_t+(H+v^2/2)_x=-a(t),\qquad 0<x<1,
\end{equation}
with
\begin{equation}\label{eq:walls}
v(t,0)=v(t,1)=0,\qquad s'(t)=a(t),\qquad D'(t)=s(t).
\end{equation}
Here $H$ is the fluid depth, $v$ its velocity relative to the tank, and $s,D$
are the tank velocity and displacement. The control $a$ is continuous, and
the fluid trajectory is jointly $C^1$ in $(t,x)$. The mathematical question
concerns classical solutions of this reduced model.

Coron proved local controllability for a sufficiently large time by the
return method \cite[Theorem 2]{Coron2002}. In his formulation of the time
problem \cite[Problem 7.1, pp.~247--250]{Coron2004}, the conjectured range
was $T>2$. Coron, Koenig and Nguyen subsequently proved a nonlinear
obstruction for $1<T<2$: an initially flat fluid cannot return to the flat
state with the same zero tank velocity under a sufficiently small nonzero
continuous control \cite[Theorem 1.2]{CKN2024}. The corresponding sufficiency
question above two is asked in AIM Problem 387 \cite{AIM387}.

We use
\[
\norm{w}_1=\max_{x\in[0,1]}(|w(x)|+|w'(x)|)
\]
and the state space
\[
\mathcal Y=\left\{(H,v,s,D)\in C^1([0,1])^2\times\mathbb R^2:
 H>0,\ \int_0^1H=1,\ v(0)=v(1)=0,\ H_x(0)=H_x(1)\right\}.
\]
Let $P(T)$ mean that for every $\varepsilon>0$ there exists $\eta>0$ such
that any two states whose total distance from $(1,0,0,0)$ is less than $\eta$
can be joined in time $T$ by a classical trajectory satisfying
\begin{equation}\label{eq:small}
\sup_{0\le t\le T}
\bigl(\norm{H(t)-1}_1+\norm{v(t)}_1+|s(t)|+|D(t)|+|a(t)|\bigr)
<\varepsilon.
\end{equation}
The conclusion below also holds with the equivalent norm
$\max|w|+\max|w'|$.

Fix throughout
\[
T=\frac{21}{10},\qquad d=\frac1{10},\qquad L=T-1=\frac{11}{10}.
\]
Define $S(r)=10r^3-15r^4+6r^5$ on $[0,1]$ and the one-periodic $C^2$
function $b$ by
\begin{equation}\label{eq:b}
b(y)=\begin{cases}
0,&0\le y\le d,\\
S((y-d)/d),&d\le y\le2d,\\
1,&2d\le y\le1-d,\\
1-S((y-(1-d))/d),&1-d\le y\le1.
\end{cases}
\end{equation}
In particular, $0\le b\le1$. For $\rho>0$ put
\begin{equation}\label{eq:target}
q_\rho(x)=\frac32\rho b'(T-x),\qquad
h_\rho(x)=\frac{q_\rho(x)+q_\rho(-x)}2,\qquad
v_\rho(x)=\frac{q_\rho(x)-q_\rho(-x)}2.
\end{equation}

\begin{theorem}\label{thm:main}
There exists $\rho_*>0$ such that, for every $0<\rho<\rho_*$, the state
$(1+h_\rho,v_\rho,0,0)$ belongs to $\mathcal Y$, and no classical trajectory
from $(1,0,0,0)$ to this state in time $T=21/10$ satisfies
\eqref{eq:small} with $\varepsilon=10^{-8}$. Consequently $P(21/10)$ is
false, and $P(T)$ does not hold for every $T>2$.
\end{theorem}

The endpoints tend to equilibrium in the prescribed norm. The theorem
therefore disproves local controllability, rather than imposing a lower bound
on the amplitude of a particular transfer. It does not give the eventual
positive controllability time.

The first-order fluid response is antiperiodic under a unit translation,
whereas the target in \eqref{eq:target} is periodic. Hence its periodic
component must come from the nonlinear terms. The tank's zero displacement
imposes a zero-mean condition on the temporal primitive producing that
component. For a time just above two, most of the quadratic spatial primitive
is the square of a single continuous function. A long positive plateau and a
short zero interval cannot be approximated by such a square while preserving
the zero mean. We use the full spatial primitive because its supremum error
can be controlled uniformly by the trajectory's $C^1$ smallness.

\section{Unfolding and uniform nonlinear estimates}
Write $h=H-1$ and unfold the fluid variables onto the period-two circle
$\Tt=\mathbb R/(2\mathbb Z)$:
\begin{equation}\label{eq:z}
z(t,x)=\begin{cases}h(t,x)+v(t,x),&0\le x\le1,\\
h(t,-x)-v(t,-x),&-1\le x\le0.
\end{cases}
\end{equation}
The wall conditions make $z$ continuous and spatially Lipschitz; it is $C^1$
on the two open halves. Reflection is understood on $\Tt$. Define
\[
\theta(x)=\begin{cases}-1,&0<x<1,\\+1,&-1<x<0,\end{cases}
\qquad
Q(A,B)=\frac{3A^2-2AB-B^2}{8},\qquad
r_z(t,x)=Q(z(t,x),z(t,-x)).
\]
Direct substitution in \eqref{eq:tank} gives
\begin{equation}\label{eq:transport}
z_t+z_x=a(t)\theta(x)-\partial_xr_z.
\end{equation}
On the positive half, $Q(z,z\circ(-\mathrm{id}))=hv+v^2/2$; on the
negative half this is the reflected flux of $h-v$. At a wall both arguments
of Q agree, so $r_z=Q(A,A)=0$. Thus the flux is continuous at each junction
and \eqref{eq:transport} holds also in distributions.

For a trajectory starting at equilibrium, let
\[
U(t)=\int_0^t a(\tau)\,d\tau,\qquad
\ell(t,x)=\int_0^t a(\tau)\theta(x+\tau-t)\,d\tau,
\]
with $U(t)=0$ for $t<0$. The linear response $\ell$ is continuous and
one-antiperiodic: $\ell(t,x+1)=-\ell(t,x)$. For $0\le x\le1$,
\begin{equation}\label{eq:linear}
\ell(t,x)=-U(t)+2\sum_{k\ge0}(-1)^kU(t-x-k).
\end{equation}
There are at most three nonzero terms for $0\le t\le T<3$. With
$A=\norm{U}_{C^0([0,T])}$ this gives $\norm{\ell}_\infty\le7A$.
Throughout this section, a supremum without a displayed domain for a
time-dependent fluid variable is over $[0,T]\times\Tt$.

Integrating \eqref{eq:transport} along straight characteristics gives
\begin{equation}\label{eq:characteristic}
z(t,x)=\ell(t,x)-\int_0^t(\partial_xr_z)(\tau,x+\tau-t)\,d\tau.
\end{equation}
There are only finitely many wall crossings; the fundamental theorem applies
on each intervening interval and the values match by continuity. Introduce
the endpoint primitive
\begin{equation}\label{eq:gamma}
\Gamma_z(x)=\int_0^T r_z(t,x+t-T)\,dt.
\end{equation}
Integrating \eqref{eq:characteristic} in x over any interval $[\xi,\zeta]$
and using continuity of the flux at its junctions gives
\[
\int_\xi^\zeta(z(T,x)-\ell(T,x))\,dx
=-\Gamma_z(\zeta)+\Gamma_z(\xi).
\]
The finite integrations can be interchanged: their bounded integrands are
continuous away from finitely many straight wall-crossing lines. Rectangular
Riemann sums agree off the cells meeting those lines, whose total area tends
to zero. Since $z(T)-\ell(T)$ is continuous, $\Gamma_z$ is $C^1$ and
\begin{equation}\label{eq:endpoint}
z(T)=\ell(T)-\partial_x\Gamma_z.
\end{equation}

\begin{lemma}\label{lem:nonlinear}
If $\sup_t(\norm{h(t)}_1+\norm{v(t)}_1)\le\varepsilon$ and
$\varepsilon\le1/(3T)$, then
\begin{equation}\label{eq:state-bound}
\norm{z}_\infty\le14A,\qquad
\norm{z-\ell}_\infty\le21T\varepsilon A\le45\varepsilon A.
\end{equation}
For pairs whose components are bounded by $M$ and $M'$ and whose
componentwise difference is at most $\delta$,
\begin{equation}\label{eq:polar}
|Q(A_1,A_2)-Q(B_1,B_2)|\le\tfrac34(M+M')\delta.
\end{equation}
\end{lemma}
\begin{proof}
On the smooth halves, $|z_x|\le\varepsilon$ and
$Q_A=(3A-B)/4$, $Q_B=-(A+B)/4$, so
$|\partial_xr_z|\le(3/2)\varepsilon\norm{z}_\infty$.
Formula \eqref{eq:characteristic} gives
$\norm{z}_\infty\le7A+(3/2)T\varepsilon\norm{z}_\infty$.
Absorbing the last term proves the first estimate; substituting it back
proves the second. To obtain \eqref{eq:polar}, integrate the gradient of Q
on the segment between the two pairs. Its componentwise sum is bounded by
$3((1-\tau)M+\tau M')/2$ at segment parameter $\tau$.
\end{proof}

\section{The quadratic endpoint and its separation}
We first compute the quadratic primitive for a linear return. For a continuous
piecewise $C^1$ function F supported in $[0,L]$ and vanishing at its endpoints,
put $U=F+F(\cdot-1)$ and $a=F'+F'(\cdot-1)$ in the linear formula. The
derivative may have finitely many jumps; this auxiliary linear control is
used only in an explicit integral. Define
\begin{equation}\label{eq:G}
J_F(s)=\int_{s-1}^{s+1}F(r)\,dr,\qquad
g_F(s)=F(s)^2-F(s)J_F(s),\qquad
G_F(s)=\sum_{n\in\mathbb Z}g_F(s+n).
\end{equation}
The sum is locally finite and $G_F$ is continuous and one-periodic.

\begin{proposition}\label{prop:quadratic}
For the linear response generated by this F, the primitive defined as in
\eqref{eq:gamma} satisfies
\begin{equation}\label{eq:quadratic}
\Gamma_\ell(x)=\frac32G_F(T-x)
 +\frac12\int_0^T U(t)^2\,dt-\frac12\int_{\mathbb R}F(t)^2\,dt.
\end{equation}
No zero-mean assumption on F is needed for this identity.
\end{proposition}
\begin{proof}
Substitution in \eqref{eq:linear} telescopes to
\begin{equation}\label{eq:paired}
\ell(t,y)=2F(t-y)-U(t),\qquad
\ell(t,-y)=U(t)-2F(t+y-1),\qquad0\le y\le1.
\end{equation}
Since $\ell$ is one-antiperiodic, its quadratic flux is one-periodic. Fix
$s=T-x$ and partition time into intervals $[\beta,\beta+1]$, where
$\beta=s+n$ for $n\in\mathbb Z$. On such an interval set
$y=t-\beta$, $P=F(\beta)$ and $R=F(2t-\beta-1)$. Expanding Q using
\eqref{eq:paired} gives
\begin{equation}\label{eq:expansion}
r_\ell(t,x+t-T)
=\frac32P^2-2PU(t)+PR+\frac12U(t)^2-\frac12R^2.
\end{equation}
Whenever $P\ne0$, the whole time interval lies in $[0,T]$, since
$0\le\beta\le L=T-1$. Also
\[
\int_\beta^{\beta+1}U(t)\,dt=J_F(\beta),\qquad
\int_\beta^{\beta+1}R\,dt=\tfrac12J_F(\beta).
\]
Thus the first three terms in \eqref{eq:expansion} integrate to
$3g_F(\beta)/2$. They vanish when $P=0$.

For the last term, the argument $r=2t-\beta-1=t+y-1$ lies in $[t-1,t]$.
Hence R is zero outside $[0,T]$ and its time integration on each interval
can be extended to that whole interval. An affine substitution gives
\[
\sum_n\int_{s+n}^{s+n+1}F(2t-s-n-1)^2\,dt
=\tfrac12\sum_n\int_{s+n-1}^{s+n+1}F(r)^2\,dr
=\int_{\mathbb R}F(r)^2\,dr.
\]
The length-two intervals at unit spacing cover each real point twice,
apart from endpoints of no effect on the integral. Only finitely many terms
intersect the support nontrivially. Summing \eqref{eq:expansion} now proves
\eqref{eq:quadratic}.
\end{proof}

The zero tank displacement will force F to have zero mean. On the long
interval $[d,1]$, this makes $G_F$ exactly $F^2$. The next lemma exploits
that relation without a derivative bound for F.

\begin{lemma}\label{lem:separation}
Let F be continuous, supported in $[0,1+d]$, with $\int_{\mathbb R}F=0$ and
$M=\norm{F}_\infty>0$. For every $c\in\mathbb R$ and $\rho>0$,
\begin{equation}\label{eq:separation}
\norm{G_F-c-\rho b}_\infty>\frac1{100}M^2.
\end{equation}
\end{lemma}
\begin{proof}
Set $I=[d,1]$, $O=[0,d]$, and for $s\in O$ write
$p(s)=F(s)$, $q(s)=F(s+1)$. On I the interval $[s-1,s+1]$ contains the
support, so $J_F=0$ and $G_F=F^2$. On O the zero mean gives
\begin{equation}\label{eq:outer}
G_F(s)=p(s)^2+q(s)^2
 +p(s)\int_s^d q(r)\,dr+q(s)\int_0^s p(r)\,dr.
\end{equation}
If $|F|$ attains M on I, then $\max G_F\ge M^2$. Otherwise the point of
maximum belongs to one of the outer intervals, and each cross term in
\eqref{eq:outer} is bounded in magnitude by $dM^2$. Thus always
$\max G_F\ge(1-2d)M^2=4M^2/5$.
If F has a zero on I, then $\min_I G_F=0$. If it has no zero there, it has
one sign; zero mean then gives $(1-d)\min_I|F|\le2dM$. Consequently
\begin{equation}\label{eq:osc}
\operatorname{osc}G_F\ge(4/5-4/81)M^2.
\end{equation}

Suppose \eqref{eq:separation} fails, and put $e=1/100$. Since b takes both
values zero and one, \eqref{eq:osc} implies
$\rho\ge(4/5-4/81-2e)M^2\ge7M^2/10$.
Let $E=\int_O(p^2+q^2)$. Integrating \eqref{eq:outer} gives
\[
\int_O G_F=E+2\int_{0\le s<r\le d}p(s)q(r)\,ds\,dr.
\]
The absolute value of the cross term is at most
$2\int_O|p|\int_O|q|\le dE$ by Cauchy--Schwarz and
$2\sqrt{ab}\le a+b$. Since $b=0$ on O,
\begin{equation}\label{eq:energy}
(1-d)E\le d(c+eM^2),\qquad c\ge-eM^2,\qquad
\left|\int_I F\right|\le\sqrt{2dE}.
\end{equation}
The last inequality uses zero mean and Cauchy--Schwarz on the two outer
intervals, whose total length is $2d$.

If $c\ge2eM^2$, then $F^2\ge c-eM^2\ge c/2$ on I. Therefore F has
one sign and $|\int_I F|\ge(1-d)\sqrt{c/2}$. But \eqref{eq:energy}
also gives $E\le3dc/[2(1-d)]$ and
$|\int_I F|\le d\sqrt{3c/(1-d)}$. This is impossible at $d=1/10$,
since $(1-d)^2/2>3d^2/(1-d)$. Hence $c<2eM^2$, so
\begin{equation}\label{eq:mean-small}
E\le\tfrac e3M^2,\qquad \left|\int_I F\right|\le\sqrt{e/15}\,M.
\end{equation}
On $[2d,1-d]$, b equals one and the presumed approximation implies
$F^2\ge\rho-2eM^2\ge17M^2/25>16M^2/25$. Thus F has one sign on
this connected interval and $|F|\ge4M/5$. Its length is $7/10$; the
remaining part of I has length $1/5$ and $|F|\le M$. It follows that
\[
\left|\int_I F\right|\ge\tfrac7{10}\tfrac45M-\tfrac15M
=\tfrac9{25}M,
\]
contradicting \eqref{eq:mean-small}.
\end{proof}

\section{Projection to a linear return and proof of the theorem}
Suppose a trajectory reaches the target \eqref{eq:target} with the smallness
bound \eqref{eq:small}. Its tank endpoints give
\begin{equation}\label{eq:moments}
U(T)=0,\qquad \int_0^T U(t)\,dt=0.
\end{equation}
We replace its linear response by one with an exact zero endpoint and the
same zero tank displacement, changing only its control primitive by a small
relative amount. This allows application of Lemma~\ref{lem:separation}.

For $0\le t\le T$ define
$F(t)=U(t)-U(t-1)+U(t-2)$, with negative-time values zero. Since $T<3$,
\[
F(t)+F(t-1)=U(t),\qquad \norm{F}_\infty\le3A.
\]
The operators $P_\pm w(x)=(w(x)\pm w(x+1))/2$ contract the supremum norm.
The unfolded terminal profile is the periodic $q_\rho$, so applying $P_-$
to \eqref{eq:state-bound} at T gives
$\norm{\ell(T)}_\infty\le45\varepsilon A$. By \eqref{eq:linear} and
\eqref{eq:moments}, $\ell(T,x)=2F(T-x)$ for $0\le x\le1$, whence
\begin{equation}\label{eq:tail}
\norm{F}_{C^0([L,T])}\le23\varepsilon A.
\end{equation}

Set $F_0(t)=F(t)-(t/L)F(L)$ on $[0,L]$, zero elsewhere. It is continuous
and piecewise $C^1$, with zero values at both support endpoints, and
$\norm{F_0-F}_{C^0([0,T])}\le23\varepsilon A$. Put
$U_0=F_0+F_0(\cdot-1)$ and $\alpha=\int_0^LF_0$.
Since $T=L+1$, $\int_0^TU_0=2\alpha$. Equation \eqref{eq:moments} gives
$|\alpha|\le23T\varepsilon A<49\varepsilon A$.

To remove this remaining mean, use
$\psi(t)=6t(L-t)/L^3$ on $[0,L]$, zero elsewhere. Its integral is one
and its supremum is $3/(2L)<3/2$. Define
\[
\widehat F=F_0-\alpha\psi,\qquad
\widehat U=\widehat F+\widehat F(\cdot-1).
\]
Then $\widehat F$ is supported in $[0,L]$, vanishes at its endpoints and
has zero mean, while
\begin{equation}\label{eq:correction}
\norm{\widehat F-F}_{C^0([0,T])}\le100\varepsilon A,\qquad
\norm{\widehat U-U}_{C^0([0,T])}\le200\varepsilon A.
\end{equation}
For $M=\norm{\widehat F}_\infty$ and $\varepsilon\le1/400$, these estimates
imply $A\le2M+200\varepsilon A$, hence $A\le4M$.

Let $\widehat\ell$ be the auxiliary linear response with primitive
$\widehat U$. It has zero terminal fluid profile by \eqref{eq:linear},
and its tank velocity and displacement endpoints vanish. Formula
\eqref{eq:linear} also gives
$\norm{\widehat\ell-\ell}_\infty\le1400\varepsilon A$. Combining
\eqref{eq:state-bound}, \eqref{eq:polar} and \eqref{eq:correction} yields
\[
\norm{z-\widehat\ell}_\infty\le1450\varepsilon A,
\qquad \norm{\widehat\ell}_\infty\le14A,
\]
and therefore
\begin{equation}\label{eq:error}
\norm{\Gamma_z-\Gamma_{\widehat\ell}}_\infty
\le T\tfrac34(28A)(1450\varepsilon A)
\le1{,}100{,}000\varepsilon M^2.
\end{equation}
The intermediate constant after using $A\le4M$ is $1{,}023{,}120$.
In particular the estimate is uniform in the frequency of the original
control, whose derivative has never been bounded.

\begin{proof}[Proof of Theorem~\ref{thm:main}]
The derivatives of S through order two match the constant pieces in
\eqref{eq:b}. Hence $q_\rho$ is $C^1$, periodic, nonzero and has zero mean.
Its even and odd parts satisfy
$\int_0^1h_\rho=0$, $v_\rho(0)=v_\rho(1)=0$ and
$h_\rho'(0)=h_\rho'(1)=0$. For all sufficiently small $\rho>0$,
$1+h_\rho>0$. These are exactly the required state conditions.

For a supposed small trajectory, apply $P_+$ to \eqref{eq:endpoint}:
\[
q_\rho=-\partial_xP_+\Gamma_z,
\qquad P_+\Gamma_z(x)=\tfrac32\rho b(T-x)+C_z
\]
for a spatial constant $C_z$. Proposition~\ref{prop:quadratic}, applied to
the corrected mean-zero $\widehat F$, gives
$\Gamma_{\widehat\ell}(x)=3G_{\widehat F}(T-x)/2+C_\ell$, which is
one-periodic. Applying the contraction $P_+$ to \eqref{eq:error}, absorbing
the constants and rotating the spatial variable gives some $c\in\mathbb R$
with
\[
\norm{G_{\widehat F}-c-\rho b}_\infty
\le\tfrac23(1{,}100{,}000)\varepsilon M^2.
\]
For $\varepsilon=10^{-8}$, all earlier smallness conditions hold and the
right-hand coefficient is $11/1500<1/100$. If $M>0$, this contradicts
Lemma~\ref{lem:separation}. If $M=0$, then $A\le4M=0$, and
\eqref{eq:state-bound} implies $z=0$, contradicting nonzero $q_\rho$.

Finally, $\norm{h_\rho}_1+\norm{v_\rho}_1=O(\rho)$ with a fixed constant.
Every endpoint neighborhood therefore contains a member of this unreachable
family, for the same fixed $T>2$ and the same fixed $\varepsilon>0$.
This disproves the asserted universal sufficiency.
\end{proof}

\paragraph{AI use.}
OpenAI Codex was used for proof development, calculation checks, retrieval of
the cited literature, and preparation and editing of this manuscript.

\begin{thebibliography}{9}
\bibitem{Coron2002}
J.-M. Coron, \emph{Local controllability of a 1-D tank containing a fluid
modeled by the shallow water equations}, ESAIM: Control, Optimisation and
Calculus of Variations \textbf{8} (2002), 513--554.
\href{https://doi.org/10.1051/cocv:2002050}{doi:10.1051/cocv:2002050}.

\bibitem{Coron2004}
J.-M. Coron, \emph{Time for local controllability of a 1-D tank containing a
fluid modeled by the shallow water equations}, Problem 7.1 in V. D. Blondel
and A. Megretski (eds.), \emph{Unsolved Problems in Mathematical Systems and
Control Theory}, Princeton University Press, 2004, 247--250.
\href{https://doi.org/10.1515/9781400826155}{doi:10.1515/9781400826155}.

\bibitem{CKN2024}
J.-M. Coron, A. Koenig and H.-M. Nguyen, \emph{Lack of local controllability
for a water-tank system when the time is not large enough}, Annales de
l'Institut Henri Poincar\'e C, Analyse non lin\'eaire \textbf{41} (2024),
1327--1365. \href{https://doi.org/10.4171/AIHPC/123}{doi:10.4171/AIHPC/123}.

\bibitem{AIM387}
AIM problem repository, \emph{387. The sharp local maneuvering time for a
shallow-water tank}, statement revision
\texttt{c9929805cece044705abaaebe1cebddac1ba11a6}.
\href{https://github.com/MColbrook/AIM/blob/c9929805cece044705abaaebe1cebddac1ba11a6/problems/387-shallow-water-tank-sharp-control-time.md}{Immutable problem statement}.
\end{thebibliography}
\end{document}
